\subsection{Conditional Expectation}
\begin{definition} \label{def:conditional-expectation-given-random-variable-elementary} \hlt{Conditional Expectation Given a Random Variable}\\
Let $X,Y$ be random variables, and assume the relevant sums or integrals exist.
\begin{enumerate}[label=(\roman*)]
\setlength{\itemsep}{0pt}
\item \hlt{Discrete case}: If $X,Y$ are discrete, then $\mathbb{E}(X\mid Y=y)\defeq\sum_x x\,p_{X\mid Y}(x\mid y)$.
\item \hlt{Density case}: If $X,Y$ have conditional density $f_{X\mid Y}$, then $\mathbb{E}(X\mid Y=y)\defeq\int_{\R}x f_{X\mid Y}(x\mid y)\,dx$
\item \hlt{Measurable transforms}: For measurable $r$,
\begin{equation*}
\mathbb{E}(r(X,Y)\mid Y=y)
=
\begin{cases}
\sum_x r(x,y)p_{X\mid Y}(x\mid y),&\text{discrete case},\\[2mm]
\int_{\R}r(x,y)f_{X\mid Y}(x\mid y)\,dx,&\text{density case}.
\end{cases}
\end{equation*}
\item \hlt{Random-variable form}: $\mathbb{E}(X\mid Y)$ is the function $g(Y)$ with $g(y)=\mathbb{E}(X\mid Y=y)$.
\end{enumerate}
\end{definition}

\begin{example} \label{ex:conditional-uniform-triangle} \hlt{A Conditional Uniform Model}\\
Let $X\sim\operatorname{Uniform}(0,1)$ and, conditionally on $X=x$, let $Y\sim\operatorname{Uniform}(x,1)$.\\
For $0\leq x<1$ and $x<y<1$, $f_{Y\mid X}(y\mid x)=1/(1-x)$, so $\mathbb{E}(Y\mid X=x)=\frac{1}{1-x}\int_x^1 y\,dy=\frac{1+x}{2}$.\\
Hence $\mathbb{E}(Y\mid X)=(1+X)/2$.
\end{example}

\begin{theorem} \label{thm:tower-property-iterated-expectation-elementary} \hlt{Tower Property}\\
Let $X,Y$ be real-valued random variables whose joint law is either discrete or admits a density, and let $r:\R^2\to\R$ be measurable.
If $\mathbb{E}|Y|<\infty$ and $\mathbb{E}|r(X,Y)|<\infty$, then $\mathbb{E}[\mathbb{E}(Y\mid X)]=\mathbb{E}Y$ and $\mathbb{E}[\mathbb{E}(r(X,Y)\mid X)]=\mathbb{E}(r(X,Y))$.
\begin{proof}
Let $h(x,y)$ denote either $y$ or $r(x,y)$, and set $g_h(x)\defeq\mathbb{E}(h(X,Y)\mid X=x)$.\\
In the density case, $g_h(x)=\int h(x,y)f_{Y\mid X}(y\mid x)\,dy$ and $f(x,y)=f_X(x)f_{Y\mid X}(y\mid x)$,\\
so Fubini gives $\mathbb{E}[g_h(X)]=\int\int h(x,y)f(x,y)\,dy\,dx=\mathbb{E}h(X,Y)$.\\
In the discrete case, $\mathbb{E}[g_h(X)]=\sum_x\sum_y h(x,y)p_{Y\mid X}(y\mid x)p_X(x)=\sum_x\sum_y h(x,y)p_{X,Y}(x,y)=\mathbb{E}h(X,Y)$.\\
Taking $h(x,y)=y$ and $h(x,y)=r(x,y)$ gives the two identities.\\
\hlt{Structural Commentary}: Conditional expectation first averages over fibers $\{X=x\}$, and the outer expectation averages those fiber means over the law of $X$.\\
\hlt{Key Insight}: Iterated expectation removes the conditioning variable without changing the total mean.
\end{proof}
\end{theorem}

\begin{example} \label{ex:tower-property-conditional-uniform} \hlt{Using the Tower Property}\\
Let $X\sim\operatorname{Uniform}(0,1)$ and, conditionally on $X=x$, let $Y\sim\operatorname{Uniform}(x,1)$.
Since $\mathbb{E}(Y\mid X)=(1+X)/2$ and $\mathbb{E}X=1/2$, the tower property gives $\mathbb{E}Y=\mathbb{E}[\mathbb{E}(Y\mid X)]=\mathbb{E}[(1+X)/2]=3/4$.
\end{example}

\begin{proposition} \label{prop:conditional-expectation-pull-out-known-factor} \hlt{Pulling Out What Is Known}\\
Let $X,Y$ be real-valued random variables whose joint law is either discrete or admits a density,\\
and let $r,s:\R\to\R$ be measurable with $\mathbb{E}|s(Y)|<\infty$, $\mathbb{E}|r(X)s(Y)|<\infty$, and $\mathbb{E}|r(X)|<\infty$.\\
Then $\mathbb{E}(r(X)s(Y)\mid X)=r(X)\mathbb{E}(s(Y)\mid X)$ and $\mathbb{E}(r(X)\mid X)=r(X)$.
\begin{proof}
Let $K_x$ be the conditional law of $Y$ given $X=x$, so $K_x(\{y\})=p_{Y\mid X}(y\mid x)$ in the discrete case and $K_x(dy)=f_{Y\mid X}(y\mid x)\,dy$ in the density case.
Set $m_s(x)\defeq\int s(y)K_x(dy)=\mathbb{E}(s(Y)\mid X=x)$, with the integral read as a sum in the discrete case.
For admissible $x$, $\mathbb{E}(r(X)s(Y)\mid X=x)=\int r(x)s(y)K_x(dy)=r(x)m_s(x)$, and $\mathbb{E}(r(X)\mid X=x)=\int r(x)K_x(dy)=r(x)$.
Evaluating at $x=X$ gives the identities.\\
\hlt{Structural Commentary}: Conditioning on $X$ makes every function of $X$ fixed inside remaining average over $Y$.\\
\hlt{Key Insight}: Known factors can be treated as constants under conditional expectation.
\end{proof}
\end{proposition}

\begin{definition} \label{def:conditional-variance-elementary} \hlt{Conditional Variance}\\
Let $X,Y$ be real-valued random variables with finite second conditional moment for each $x$ under consideration.
\begin{enumerate}[label=(\roman*)]
\setlength{\itemsep}{0pt}
\item The conditional mean is $\mu(x)\defeq\mathbb{E}(Y\mid X=x)$.
\item The conditional variance at $x$ is $v(x)\defeq\operatorname{Var}(Y\mid X=x)\defeq\mathbb{E}((Y-\mu(x))^2\mid X=x)$.
\item The conditional-variance random variable is $\operatorname{Var}(Y\mid X)\defeq v(X)$.
\end{enumerate}
\end{definition}

\begin{theorem} \label{thm:law-total-variance} \hlt{Law of Total Variance}\\
Let $X,Y$ be real-valued random variables with $\mathbb{E}Y^2<\infty$.
Set $b(x)\defeq\mathbb{E}(Y\mid X=x)$ and $v(x)\defeq\operatorname{Var}(Y\mid X=x)$, and write $b(X)$ and $v(X)$ for the induced random variables.
Then $\operatorname{Var}(Y)=\mathbb{E}[v(X)]+\operatorname{Var}(b(X))$.
\begin{proof}
Let $m\defeq\mathbb{E}Y$ and $e\defeq Y-b(X)$.\\
The tower property gives $\mathbb{E}b(X)=m$, and $\mathbb{E}(e\mid X)=\mathbb{E}(Y\mid X)-b(X)=0$.\\
Since $Y-m=e+(b(X)-m)$, expansion gives $\operatorname{Var}(Y)=\mathbb{E}e^2+2\mathbb{E}[e(b(X)-m)]+\mathbb{E}(b(X)-m)^2$.\\
The cross term is zero because $\mathbb{E}[e(b(X)-m)]=\mathbb{E}[(b(X)-m)\mathbb{E}(e\mid X)]=0$.\\
Also $\mathbb{E}e^2=\mathbb{E}[\mathbb{E}(e^2\mid X)]=\mathbb{E}[v(X)]$, $\mathbb{E}(b(X)-m)^2=\operatorname{Var}(b(X))$.\\
By substitution, $\operatorname{Var}(Y)=\mathbb{E}[v(X)]+\operatorname{Var}(b(X))$.\\
\hlt{Structural Commentary}: The decomposition separates within-fiber variation $v(X)$ from between-fiber variation of the conditional mean $b(X)$.\\
\hlt{Key Insight}: Total variance equals average conditional variance plus variance of conditional expectation.
\end{proof}
\end{theorem}

\begin{example} \label{ex:total-variance-random-binomial-proportion} \hlt{Random Binomial Proportion}\\
Let $Q\sim\operatorname{Uniform}(0,1)$ and, conditionally on $Q=q$, let $X\sim\operatorname{Binomial}(n,q)$.\\
Then $\mathbb{E}(X\mid Q)=nQ$, so $\mathbb{E}X=\mathbb{E}(nQ)=n/2$. Also $\operatorname{Var}(X\mid Q)=nQ(1-Q)$.\\
Hence $\mathbb{E}[\operatorname{Var}(X\mid Q)]=n\int_0^1q(1-q)\,dq=n/6$, $\operatorname{Var}(\mathbb{E}(X\mid Q))=\operatorname{Var}(nQ)=n^2/12$.\\
Therefore $\operatorname{Var}(X)=n/6+n^2/12$.
\end{example}

\begin{corollary} \label{cor:constant-conditional-mean-implies-uncorrelated} \hlt{Constant Conditional Mean Implies Uncorrelated}\\
Let $X,Y\in L^2$ and suppose $\mathbb{E}(X\mid Y)=c$ a.s. for some constant $c\in\R$. Then $c=\mathbb{E}X$ and $\operatorname{Cov}(X,Y)=0$.
\begin{proof}
The tower property gives $\mathbb{E}X=\mathbb{E}[\mathbb{E}(X\mid Y)]=c$.\\
Since $Y$ is known when conditioning on $Y$, $\mathbb{E}(XY\mid Y)=Y\mathbb{E}(X\mid Y)=cY$.\\
Taking expectations gives $\mathbb{E}(XY)=c\mathbb{E}Y=\mathbb{E}X\,\mathbb{E}Y$, hence $\operatorname{Cov}(X,Y)=\mathbb{E}(XY)-\mathbb{E}X\,\mathbb{E}Y=0$.\\
\hlt{Structural Commentary}: A constant conditional mean removes all linear dependence detectable through averaging against functions of $Y$; in particular it removes covariance with $Y$.\\
\hlt{Key Insight}: Mean independence implies uncorrelatedness, but not conversely.
\end{proof}
\end{corollary}
